В прошлых разделах мы научились соединять базовые алгоритмы в ансамбль с помощью bagging
(и, в частности, строить из решающих деревьев случайные леса).
Теперь мы рассмотрим другой способ объединять базовые алгоритмы в композицию -- градиентный бустинг.

В ходе обучения случайного леса каждый базовый алгоритм строится независимо от остальных.
Бустинг, в свою очередь, воплощает идею последовательного построения линейной комбинации алгоритмов.
Каждый следующий алгоритм старается уменьшить ошибку текущего ансамбля.

Бустинг, использующий деревья решений в качестве базовых алгоритмов, называется градиентным бустингом над решающими
деревьями, (Gradient Boosting on Decision Trees, GBDT).

Он отлично работает на выборках с «табличными», неоднородными данными.
Пример таких данных -- описание пользователя Яндекса через его возраст, пол, среднее число поисковых запросов в
день, число заказов такси и так далее.
Такой бустинг способен эффективно находить нелинейные зависимости в данных различной природы.

Этим свойством обладают все алгоритмы, которые используют деревья решений, однако именно GBDT обычно выигрывает в
подавляющем большинстве задач.
Благодаря этому он широко применяется во многих конкурсах по машинному обучению и задачах из индустрии

\begin{itemize}
    \item В поисковом ранжировании;
    \item рекомендательных системах;
    \item таргетировании рекламы;
    \item предсказании погоды;
    \item выбора пункта назначения такси и многих других.
\end{itemize}

Не так хорошо бустинг проявляет себя на однородных данных: текстах, изображениях, звуке, видео.
В таких задачах нейросетевые подходы почти всегда демонстрируют лучшее качество.

И хотя деревья решений -- традиционный выбор для объединения в ансамбли, никто не запрещает использовать и другие
алгоритмы (например, линейные модели) в качестве базовых.
Эта возможность реализована в пакете XGBoost.

Стоит только понимать, что построенная композиция окажется линейной комбинацией линейных моделей, то есть опять-таки
линейной моделью - или нейросетью с одним полносвязным слоем.
Это уменьшает возможности ансамбля эффективно определять нелинейные зависимости в данных.
Поэтому в этом параграфе мы рассмотрим только бустинг над решающими деревьями.

\subsection{Процедура построения для MSE}

Рассмотрим задачу регрессии с квадратичной функцией потерь:
\[ L(y, x) = \frac{1}{2} \sum_{i=1}^{N} \bigl(y_i - a(x_i)\bigr)^2 \; \to \; \min \]

Для решения будем строить композицию из $K$ базовых алгоритмов:\[ a(x) = a_K(x) = b_1(x) + b_2(x) + \dots + b_K(x) \]

Если мы обучим единственное решающее дерево, то качество такой модели, скорее всего, будет низким.
Однако мы знаем, на каких объектах построенное дерево давало точные предсказания, а на каких ошибалось.

Попробуем использовать эту информацию и обучим ещё одну модель.
Допустим, что предсказание первой модели на $x_l$ на 10 больше, чем необходимо (то есть $b_1(x_l) = y_l + 10$).
Если бы мы могли обучить новую модель, которая на $x_l$ будет выдавать ответ $-10$, то сумма ответов этих двух
моделей на объекте $x_l$ в точности совпала бы с истинным значением:
\[ b_1(x_l) + b_2(x_l) = (y_l + 10) + (-10) = y_l \]

Другими словами, если вторая модель научится предсказывать разницу между реальным значением и ответом первой, то это
позволит уменьшить ошибку композиции.

В реальности вторая модель тоже не сможет обучиться идеально, поэтому обучим третью, которая будет «компенсировать»
неточности первых двух.
Будем продолжать так, пока не построим композицию из $K$ алгоритмов.

Для объяснения метода градиентного бустинга полезно воспользоваться следующей аналогией.
Бустинг можно представить как гольфиста, цель которого -- загнать мяч в лунку с координатой $y_{\text{ball}}$.
Положение мяча здесь -- ответ композиции $a(x_{\text{ball}})$.
Гольфист мог бы один раз ударить по мячу, не попасть в лунку и пойти домой, но настырность заставляет его продолжить.

По счастью, ему не нужно начинать каждый раз с начальной позиции.
Следующий удар переводит мяч из текущего положения $a_k(x_{\text{ball}})$ в положение $a_{k+1}(x_{\text{ball}})$.
Каждый следующий удар -- это та поправка, которую вносит очередной базовый алгоритм в композицию.
Если гольфист всё делает правильно, то функция потерь будет уменьшаться:
\[ L\bigl(y, a_{k+1}(x)\bigr) < L\bigl(y, a_k(x)\bigr),\]

То есть мяч постепенно будет приближаться к лунке.

Удары при этом делаются не хаотично.
Гольфист оценивает текущее положение мяча относительно лунки и следующим ударом старается нивелировать те проблемы,
которые он создал себе всеми предыдущими.
Подбираясь к лунке, он будет бить всё аккуратнее и, возможно, даже возьмёт другую клюшку, но точно не будет лупить
так же, как из первоначальной позиции.
В итоге комбинация всех ударов рано или поздно перенесет мяч в лунку.

Подобно тому, как гольфист постепенно подводит мяч к цели, бустинг с каждым новым базовым алгоритмом всё больше
приближает предсказание к истинному значению метки объекта.

Рассмотрим теперь другую аналогию — разложение функции в ряд Тейлора.
Из курса математического анализа известно, что (достаточно хорошую) бесконечно дифференцируемую функцию $f(x)$ на
интервале $x \in (a-R, a+R)$ можно представить в виде бесконечной суммы степенных функций:
\[ f(x) = \sum_{n=0}^{\infty} \frac{f^{(n)}(a)}{n!} (x-a)^n. \]

Одна, самая первая степенная функция в разложении, очень грубо приближает $f(x)$.
Прибавляя к ней следующую, мы получим более точное приближение.
Каждая следующая элементарная функция увеличивает точность приближения, но менее заметна в общей сумме.
Если нам не требуется абсолютно точное разложение, вместо бесконечного ряда Тейлора мы можем ограничиться суммой его
первых $k$ элементов.
Таким образом, интересующую нас функцию мы с некоторой точностью представили в виде суммы «простых» функций.

Перенесём эту идею на задачи машинного обучения.
В машинном обучении мы пытаемся по выборке $(x_i,y_i)$ восстановить неизвестную истинную зависимость.
Прежде всего, мы выбираем подходящий алгоритм.
Мы можем выбрать «сложный» алгоритм, который сразу хорошо выучит истинную зависимость.

А можем обучить «простой», который выучит истинную зависимость посредственно.
Затем мы добавим к нему ещё один такой простой алгоритм, чтобы уточнить предсказание первого алгоритма.
Продолжая этот процесс, мы получим сумму простых алгоритмов, где первый алгоритм грубо приближает истинную
зависимость, а каждый следующий делает приближение всё точнее.

\subsection{Бустинг как градиентный спуск}

Градиентный бустинг (Gradient Boosting) можно рассматривать как градиентный спуск в пространстве функций.
Вместо оптимизации параметров $\bm{\theta}$ модели $f(\bm{x}; \bm{\theta})$, мы оптимизируем саму функцию
$f: \mathbb{R}^d \to \mathbb{R}$.

\underline{Формальная постановка}

Рассмотрим задачу минимизации эмпирического риска:

\[ \min_{f \in \mathcal{F}} \mathcal{L}(f) = \sum_{i=1}^n L(y_i, f(\bm{x}_i)) \]

где:
\begin{itemize}
    \item $L: \mathbb{R} \times \mathbb{R} \to \mathbb{R}_+$ --- функция потерь
    \item $\mathcal{F}$ --- пространство функций (например, деревья решений)
    \item $f(\bm{x})$ --- наша модель
\end{itemize}

\underline{Градиентный спуск в пространстве функций}

На каждой итерации $t$ обновляем функцию:
\[ f_t(\bm{x}) = f_{t-1}(\bm{x}) - \eta_t \cdot g_t(\bm{x}) \]

Здесь $g_t(\bm{x})$ --- \("\)градиент\("\) функции потерь по $f$ в точке $f_{t-1}$.

\underline{Градиент функции потерь}

Для фиксированной точки $(\bm{x}_i, y_i)$ рассмотрим производную:

\[ g_i = \left. \frac{\partial L(y_i, z)}{\partial z} \right|_{z = f_{t-1}(\bm{x}_i)} \]

Нам нужно найти функцию $h_t \in \mathcal{F}$, которая лучше всего аппроксимирует отрицательный градиент:

\[ h_t \approx -g_t \]

Здесь $g_t(\bm{x}_i) = g_i$

\subsubsection{вывод для кросс-энтропии}

\underline{Функция потерь и вероятностная модель}

Для бинарной классификации с $y_i \in \{0, 1\}$ используем кросс-энтропийную потерь:

\[ L(y, f(\bm{x})) = -y \log(p(\bm{x})) - (1-y) \log(1 - p(\bm{x})) \]

Здесь $p(\bm{x}) = P(y=1 \mid \bm{x}) = \sigma(f(\bm{x})) = \frac{1}{1 + e^{-f(\bm{x})}}$.

\underline{Вычисление градиента}

Выразим потерю через $f(\bm{x})$:

\begin{align*}
    L(y, f) &= -y \log\left(\frac{1}{1 + e^{-f}}\right) - (1-y) \log\left(1 - \frac{1}{1 + e^{-f}}\right) \\
    &= -y \log\left(\frac{1}{1 + e^{-f}}\right) - (1-y) \log\left(\frac{e^{-f}}{1 + e^{-f}}\right) \\
    &= y \log(1 + e^{-f}) + (1-y) \left[\log(1 + e^{-f}) - \log(e^{-f})\right] \\
    &= y \log(1 + e^{-f}) + (1-y) \left[\log(1 + e^{-f}) + f\right] \\
    &= \log(1 + e^{-f}) + (1-y)f
\end{align*}

Теперь вычислим градиент:

\begin{align*}
    \frac{\partial L(y, f)}{\partial f} &= \frac{\partial}{\partial f} \left[ \log(1 + e^{-f}) + (1-y)f \right] \\
    &= \frac{-e^{-f}}{1 + e^{-f}} + (1-y) \\
    &= -\frac{1}{1 + e^{f}} + (1-y) \\
    &= (1-y) - \sigma(-f) \\
    &= (1-y) - \frac{1}{1 + e^{f}}
\end{align*}

Упростим дальше.
Заметим, что $\sigma(-f) = 1 - \sigma(f)$, поэтому:

\begin{align*}
    \frac{\partial L(y, f)}{\partial f} &= (1-y) - [1 - \sigma(f)] \\
    &= \sigma(f) - y
\end{align*}

Таким образом, получаем простую формулу:

\[ \boxed{\frac{\partial L(y, f)}{\partial f} = p(\bm{x}) - y} \]

Здесь $p(\bm{x}) = \sigma(f(\bm{x}))$.

\underline{Интерпретация градиента}

Для каждого объекта $(\bm{x}_i, y_i)$:

\[ g_i = \left. \frac{\partial L(y_i, f)}{\partial f} \right|_{f = f_{t-1}(\bm{x}_i)} = p_{t-1}(\bm{x}_i) - y_i \]

Здесь $p_{t-1}(\bm{x}_i) = \sigma(f_{t-1}(\bm{x}_i))$ --- предсказанная вероятность класса 1 на предыдущей итерации.

\begin{itemize}
    \item Если $y_i = 1$ и $p_{t-1}(\bm{x}_i) = 0.9$, то $g_i = -0.1$ (немного уменьшаем $f$)
    \item Если $y_i = 0$ и $p_{t-1}(\bm{x}_i) = 0.9$, то $g_i = 0.9$ (сильно увеличиваем $f$)
    \item Если $y_i = 1$ и $p_{t-1}(\bm{x}_i) = 0.1$, то $g_i = -0.9$ (сильно уменьшаем $f$)
\end{itemize}

\underline{Практические аспекты}

\underline{Регуляризация}

\begin{enumerate}
    \item \textbf{Скорость обучения (learning rate)} $\eta$: значения ($\eta < 0.1$) улучшают обобщающую способность
    \item \textbf{Глубина деревьев}: ограничивает сложность базовых моделей
    \item \textbf{Subsampling}: использование подвыборки объектов на каждой итерации
    \item \textbf{Column subsampling}: использование подвыборки признаков
\end{enumerate}

\underline{Преимущества градиентной интерпретации}

\begin{enumerate}
    \item \textbf{Гибкость}: можно использовать любую дифференцируемую функцию потерь
    \item \textbf{Регуляризация}: легко добавлять регуляризаторы
    \item \textbf{Интерпретируемость}: можем анализировать вклады каждой итерации
\end{enumerate}

Таким образом
\begin{itemize}
    \item Градиентный бустинг можно рассматривать как градиентный спуск в пространстве функций
    \item Для кросс-энтропийной потери градиент имеет простую форму: $p(\bm{x}) - y$
    \item Это объясняет, почему на каждой итерации мы обучаем модель на остатках $y - p(\bm{x})$
    \item Градиентная интерпретация позволяет обобщать бустинг на различные функции потерь и задачи
\end{itemize}

Градиентный бустинг остаётся одним из самых мощных методов машинного обучения, и понимание его как градиентного спуска
помогает как в его применении, так и в разработке новых модификаций.

\subsection{Выбор базовой модели}

Why not linear?
Why not kNN or others?
Why trees as basic algorithms?
\begin{enumerate}
    \item GB doesn’t need differentiable base model
    \subitem target is given independently between iterations
    \item Trees have high variance
    \subitem gradient varies a lot between iterations
    \subitem это означает высокую дисперсию целевых показателей для базовой модели
    \item Error on any step doesn’t ruin whole algorithm
    \subitem we recompute true gradients on each step
    \subitem meta-algorithm (GB) fixes such errors
\end{enumerate}

So for now trees are the only option for GB

\subsection{Параллелизм}

Сам градиентный бустинг не parallelized, но так как обычно его делают на деревьях, то можно распараллелить процесс
построения каждого дерева